% Данный файл распространяетсяя под лицензией CC BY 4.0. Текст лицензии размещён на https://creativecommons.org/licenses/by/4.0/
% (c) Кафедра системного программирования, 2025

\documentclass[a4paper]{article}

\usepackage[a4paper, top=8mm, bottom=8mm, left=8mm, right=8mm]{geometry}

\usepackage{polyglossia}
\setdefaultlanguage[babelshorthands=true]{russian}
\setotherlanguage{english}

\usepackage{fontspec}
\setmainfont{FreeSerif}
\newfontfamily{\russianfonttt}[Scale=0.7]{DejaVuSansMono}

\usepackage[tiny, compact]{titlesec}

\usepackage{titling}
\setlength{\droptitle}{-1cm}
\pretitle{\begin{center}\begin{bfseries}\Large}
\posttitle{\par\end{bfseries}\end{center}}
\preauthor{\begin{center}\normalsize}
\postauthor{\par\end{center}\vspace{-1.8cm}}

\usepackage{hyperref}
\usepackage{bookmark}
\usepackage{csquotes}

\title{Добавление алгебраических типов данных в компилятор rukaml}

\author{Балышев Артем Максимович}

\date{}

\begin{document}

\maketitle

\begin{flushright}
    Группа: \emph{\enquote{24.Б72-мм}}

    Кафедра: \emph{Системного программирования}

    Научный руководитель: \emph{Косарев Дмитрий Сергеевич}

    Консультант: \emph{Косарев Дмитрий Сергеевич}

    Номер семестра практики: \emph{3}
\end{flushright}


\section{Постановка задачи}

Работа выполняется в рамках проекта \enquote{rukaml} по разработке компилятора языка OCaml. \textbf{Целью работы является} реализация поддержки алгебраических типов данных и паттерн-матчинга выражений.

Конечным пользователем продукта являются другие разработчики, в частности компиляторные инженеры. Алгебраические типы данных позволяют представлять программу на различных этапах компиляции в виде, удобном для проверки её корректности и дальнейшей кодогенерации или интерпретации.

\section{Описание предлагаемого решения}

В рамках работы предлагается добавление возможности объявления пользовательских типов и проверка синтаксической и семантической корректности объявлений типов и использований конструкторов типов.
Для работы с конструкторами алгебраических типов реализуется поддержка паттерн-матчинга с дальнейшим преобразованием его в деревья выбора (англ. Decision Trees).
Также предлагается добавление списков как частного случая алгебраических типов данных в качестве встроенного типа данных.


Реализация решения осуществлялась на языке OCaml с использованием системы сборки dune и некоторых вспомогательных утилит и библиотек: Angstrom, Jane Street Base, bisect\_ppx, ppx\_deriving.

\section{Тестирование}

В проекте rukaml для тестирования реализуются отдельные модули, называемые pretty-printer'ами, которые позволяют распечатать представление программы, полученное после её обработки компилятором, для последующей ручной проверки соответствия реального вывода компилятора ожидаемому.
Для обработки ошибок синтаксического анализатора и семантического анализатора используются монадические интерфейсы, предоставляемые библиотекой Angstrom и стандартной библиотекой языка OCaml.
Тестирование проводится с помощью встроенной в систему сборки dune поддержки cram-тестов — текстовых сценариев командной оболочки.

\section{Заключение}

В ходе выполнения работы реализована следующая функциональность:

\begin{itemize}
    \item Построение деревьев абстрактного синтаксиса для объявлений алгебраических типов данных и выражений с паттерн-матчингом
    \item Вывод типов выражений c паттерн-матчингом и выражений, являющихся экземплярами алгебраических типов данных
    \item Поддержка списков в качестве встроенного типа данных на этапах синтаксического разбора и семантического разбора
    \item Преобразование паттерн-матчинга выражений в деревья выбора
\end{itemize}

Репозиторий проекта: \url{https://github.com/Kakadu/rukaml}


Ссылки на pull requests:
\begin{itemize}
    \item \url{https://github.com/Kakadu/rukaml/pull/12}
    \item \url{https://github.com/Kakadu/rukaml/pull/18}
    \item \url{https://github.com/Kakadu/rukaml/pull/19}
\end{itemize}


Техническая документация: \url{https://github.com/psiblvdegod/third-sem-practice/wiki/Документация}


Работа находится в ветке dev проекта, в будущем будет объединена с кодом в ветке master.

\end{document}
